\documentclass[10pt,a4paper]{amsart}
\usepackage[margin=19mm]{geometry}
\usepackage{amsmath,amssymb,mathtools}
\usepackage[scr=rsfs,cal=euler]{mathalfa}
\usepackage{xfrac}
\usepackage[unicode,hidelinks,bookmarks=false]{hyperref}
\newcommand{\ie}{i.e.\ }
\newcommand{\cf}{cf.\ }
\newcommand{\resp}{respectively\ }
\newcommand{\Subsection}{\subsection}
\providecommand{\nobreakdash}{\nobreak\hbox{-}}
\setlength{\emergencystretch}{2em}
\multlinegap0pt
\usepackage{amssymb}
\usepackage{mathtools}
\usepackage[shortlabels]{enumitem}
\usepackage{xcolor}
\newtheorem{proposition}{Proposition}
\newcommand{\C}{\mathbb{C}}
\newcommand{\sigthree}{\sigma_3}
\title{The Domb Apéry-limit and a proof of the Ramanujan Machine conjecture Z2}
\author{Alex Shvets}
\date{Source: arXiv:2604.06239v1 (CC BY 4.0)}
\begin{document}
\maketitle

\noindent\emph{Source and licence.} The Domb Apéry-limit and a proof of the Ramanujan Machine conjecture Z2. Version arXiv:2604.06239v1 (CC BY 4.0), distributed by arXiv under the Creative Commons Attribution 4.0 International licence, http://creativecommons.org/licenses/by/4.0/, as recorded in the arXiv licence metadata for this exact version and shown on the abstract page https://arxiv.org/abs/2604.06239v1. Full licence notice: \texttt{licenses/arxiv-2604.06239-CC-BY-4.0.txt}.\par\smallskip

\noindent\textit{Editorial scope.} Counted unit: the article's proposition prop:Lfunctions (source0.tex lines 372--403). The article's complete original proof of it is delivered with it (source0.tex lines 405--442, one \texttt{proof} environment of 38 lines). No mathematics is written, repaired or re-derived: the statement and the proof are the article's own lines, in the article's own notation, and no retained line is substituted or re-flowed.  Retained uncounted as the source-local context the proof needs: lines 340--342.  Nothing was omitted from any retained span, so the package carries no omission record.  No retained line contains a citation, so the wrapper carries no bibliography and no bibliography entry.  No retained line contains a figure environment, an image-inclusion command or drawing code, so no image is delivered and the package declares no asset dependency.  Editorial formatting only, in addition: an AMS \texttt{amsart} preamble that restates the article's own declarations and macros used by the retained lines, namely the environments \texttt{proposition} on separate counters, together with the standard packages those lines need.  The retained environments print in this document as Proposition 1 in order of appearance, whereas the article prints the same items with the numbering of its own full document.  The complete native source of the article as distributed in the arXiv e-print for this exact version is bundled unchanged as \texttt{sources/198009/source0.tex}.
\begin{equation}\label{eq:a-n-def}
a_n=\sigthree(n)-\sigthree(n/2)-9\sigthree(n/3)+9\sigthree(n/6).
\end{equation}

\begin{proposition}\label{prop:Lfunctions}
The Dirichlet series of $g$ is
\begin{equation}\label{eq:Lg-factor}
L(g,s):=\sum_{n\ge 1}\frac{a_n}{n^s}=\zeta(s)\zeta(s-3)(1-2^{-s})(1-3^{2-s}).
\end{equation}
Hence
\begin{equation}\label{eq:Lg3}
L(g,3)=-\frac{7}{24}\zeta(3).
\end{equation}
Define the alternating twist
\begin{equation}\label{eq:Lstar-def}
L^*(s):=\sum_{n\ge 1}\frac{(-1)^n a_n}{n^s}.
\end{equation}
Then
\begin{equation}\label{eq:Lstar-factor}
L^*(s)=-\zeta(s)\zeta(s-3)(1-2^{-s})(1-2^{4-s})(1-3^{2-s}).
\end{equation}
In particular,
\begin{equation}\label{eq:Lstar-specials}
L^*(3)=-\frac{7}{24}\zeta(3),
\qquad
L^*(2)=0,
\qquad
L^*(1)=\frac{7}{4\pi^2}\zeta(3).
\end{equation}
Moreover, for every $s\in\C$,
\begin{equation}\label{eq:Lstar-fe}
\Gamma(s)\left(\frac{\pi}{\sqrt 3}\right)^{-s}L^*(s+3)
=
\Gamma(-s-2)\left(\frac{\pi}{\sqrt 3}\right)^{s+2}L^*(1-s).
\end{equation}
\end{proposition}

\begin{proof}
Equation \eqref{eq:Lg-factor} follows immediately from \eqref{eq:a-n-def} and the classical identity
\[
\sum_{n\ge 1}\frac{\sigthree(n)}{n^s}=\zeta(s)\zeta(s-3).
\]
Substituting $s=3$ yields \eqref{eq:Lg3} because $\zeta(0)=-1/2$.

For the twist, write $n=2^rm$ with $m$ odd. Since \eqref{eq:a-n-def} implies
\[
a_{2^rm}=8^r a_m\qquad(m\text{ odd},\ r\ge 0),
\]
we get
\begin{align*}
L^*(s)
&=\sum_{m\text{ odd}}\frac{a_m}{m^s}\left(-1+\sum_{r\ge 1}\frac{8^r}{2^{rs}}\right)\\
&=\sum_{m\text{ odd}}\frac{a_m}{m^s}\left(-1+\frac{2^{3-s}}{1-2^{3-s}}\right)
=-\frac{1-2^{4-s}}{1-2^{3-s}}\sum_{m\text{ odd}}\frac{a_m}{m^s}.
\end{align*}
On the other hand, removing the $2$-Euler factor from \eqref{eq:Lg-factor} gives
\[
\sum_{m\text{ odd}}\frac{a_m}{m^s}=(1-2^{3-s})L(g,s).
\]
Combining these formulas with \eqref{eq:Lg-factor} proves \eqref{eq:Lstar-factor}. The values in \eqref{eq:Lstar-specials} follow immediately, using the standard identity
\[
\zeta'(-2)=-\frac{\zeta(3)}{4\pi^2}.
\]

Finally, \eqref{eq:Lstar-fe} is a direct consequence of \eqref{eq:Lstar-factor} and the functional equation of the Riemann zeta function. Indeed,
\begin{align*}
&\Gamma(s)\left(\frac{\pi}{\sqrt 3}\right)^{-s}L^*(s+3)\\
&\qquad =-\Gamma(s)\left(\frac{\pi}{\sqrt 3}\right)^{-s}\zeta(s+3)\zeta(s)(1-2^{-s-3})(1-2^{1-s})(1-3^{-s-1})
\end{align*}
and the two zeta functional equations transform the right-hand side into
\[
\Gamma(-s-2)\left(\frac{\pi}{\sqrt 3}\right)^{s+2}L^*(1-s).
\qedhere
\]
\end{proof}

\end{document}
